%======================================================================
\subsection{UI design mock: clean but feature dense (19 Oct 2026)}\label{sec:ui-design}
%======================================================================

\textbf{Your answers (19 Oct).}
\begin{itemize}[nosep]
  \item A design mock first, then this plan, approved before any code.
  \item Light and dark, following the macOS appearance setting; the accent follows the macOS accent colour.
  \item ``Clean but feature dense.''
  \item Keep FreeCAD's icons, but at one size with consistent captions.
  \item BruceWorx's own top bar (workspace menu, tabs, captioned groups) replaces the FreeCAD-Ribbon add-on.
  \item Screens: part modelling, assembly with the joint panel, sketch, and the Start page.
\end{itemize}

The mock is HTML drawn at the size of your screen (1600$\times$1000 points, 2$\times$). It uses FreeCAD's own icons,
and its 3D views are rendered by this build. It shows what is \emph{intended}. It is not a screenshot of working
code. Every element is mapped below to how it would be built. The source is in
\file{scripts/ui\_mock/}, and \code{render.sh} redraws every figure in this section.

\subsubsection*{Where the window is today}

\begin{figure}[!htbp]
\centering
\begin{tikzpicture}
  \node[anchor=south west,inner sep=0] (img) at (0,0) {\includegraphics[width=\linewidth]{fig/ui_before_joint.png}};
  \begin{scope}[x={(img.south east)},y={(img.north west)}]
    \foreach \n/\x/\y in {1/0.035/0.972, 2/0.435/0.972, 3/0.985/0.977, 4/0.10/0.78, 5/0.175/0.53, 6/0.78/0.79, 7/0.80/0.49, 8/0.12/0.145, 9/0.05/0.105, 10/0.61/0.026}
      \node[circle,fill=bwred,text=white,font=\scriptsize\bfseries,inner sep=1.2pt,minimum size=13pt,draw=white,line width=0.8pt] at (\x,\y) {\n};
  \end{scope}
\end{tikzpicture}
\caption{This build tonight (the hinge from tests 14 and 16), with the problems numbered.}
\label{fig:ui-before}
\end{figure}

\begin{center}
\small
\begin{tabularx}{\linewidth}{@{}cXX@{}}
\toprule
\# & What you see & Cause and fix\\
\midrule
\marker{1} & ``Menu'' button, 13 workbench tabs, a large icon grid & The \textbf{FreeCAD-Ribbon} add-on (v1.12.0dev) in your FreeCAD profile, not BruceWorx. It is replaced by the BruceWorx bar; see question~3.\\
\marker{2} & Title shows ``FreeCAD 1.1.4 -- Hinge : 1*'' & Window title: the document name only, with the unsaved mark on its tab.\\
\marker{3} & Windows-style $-\,\square\,\times$ at the top right & Probably drawn by Ribbon. \verify{} Native macOS buttons (\eref{UIX-02}) come back once it is off.\\
\marker{4} & Model tree names cut short and faded (``Ass\ldots bly'', ``F\ldots e'') & The faded names have about 1.1:1 contrast. Suspected: the floating-tree sheet's white text (\file{BruceWorxModelTree.qss}) applied to a light panel, and a narrow first column. \verify\\
\marker{5} & Property headers grey on grey & Header colour from the Ribbon sheet; replaced by the token colours below.\\
\marker{6} & ``Tasks'' column the full window height, header floating & Task panel docked right. It becomes a floating command dialog with OK/Cancel at the bottom.\\
\marker{7} & Odd gaps, a stray icon, one \emph{Clear} button styled differently & Task panel layout and QSS (\eref{UIX-07}).\\
\marker{8} & Document tabs at the bottom, with red $\times$ & \code{setTabPosition(South)} is hard-coded at \file{MainWindow.cpp:397}; the red $\times$ is a stylesheet icon. \verify\\
\marker{9} & Timeline with two headers that says ``No PartDesign body'' in an assembly & Dock title plus bar header. It becomes one row and is hidden when there is no history.\\
\marker{10} & Status bar: black bars, ``Revit'', a size cut off & Status-bar widgets with no styling. \verify{} They are regrouped as quiet readouts.\\
\bottomrule
\end{tabularx}
\end{center}

\subsubsection*{The design, one screen explained}

\begin{figure}[!htbp]
\centering
\begin{tikzpicture}
  \node[anchor=south west,inner sep=0] (img) at (0,0) {\includegraphics[width=\linewidth]{fig/ui_mock_part-notes.png}};
\end{tikzpicture}
\caption{Part modelling, light. A body with a fillet being edited. The numbers match the table below.}
\label{fig:ui-notes}
\end{figure}

\begin{center}
\small
\begin{tabularx}{\linewidth}{@{}cp{3.3cm}Xl@{}}
\toprule
\# & Element & How it is built & Epic\\
\midrule
\marker{1} & Document tabs in the title bar & Tab bar moved \code{North} (a preference instead of the hard-coded \code{South}), then hosted in the unified title bar. Modified dot instead of $\times$; $\times$ on hover. & \eref{UIX-06}\\
\marker{2} & Search, \texttt{Cmd-K} & Commands, preferences and tree in one field. Built in, or the SearchBar add-on restyled (question~2). & \eref{UIX-06}\\
\marker{3} & Vault status & Small Python status widget reading the folder vault (\S\ref{sec:pdm-eval}). & \eref{PDM-04}\\
\marker{4} & Workspace menu & \code{WorkspaceSelector}, as in the starting brief. & \eref{UIX-05}\\
\marker{5} & Tabs, plus \textsc{Sketch} in accent while a sketch is open & Toolbar sets per tab; the contextual tab is pushed on sketch edit. & \eref{UIX-05}\\
\marker{6} & Captioned groups, 24\,px icons & One icon size (preference), each toolbar given a caption and a $\blacktriangledown$ overflow menu. No icon is redrawn. & \eref{UIX-06}\\
\marker{7} & Browser with properties & Tree and property view in the left overlay as one card. Second column shows the feature summary (``Pad 8 mm''), P/N and state badges from the part properties (\S\ref{sec:partprops}). & \eref{UIX-03}\\
\marker{8} & Command dialog & Task panel in the right overlay as a card. Icon and title header, OK/Cancel at the bottom (\code{TaskDialog::South} as default), keyboard hint. & \eref{UIX-07}\\
\marker{9} & View cube & NaviCube size, colours and font are already preferences. & \eref{UIX-01}\\
\marker{10} & Navigation pill & Small Python overlay on the 3D view: orbit, pan, zoom, fit, display style, camera, section, measure. & \eref{UIX-06}\\
\marker{11} & Timeline, one row & HistoryTimeline mod: no dock title, 20\,px icons, rolled-back items dimmed. & \eref{UIX-04}\\
\marker{12} & Status bar & Prompt on the left; Report, units, navigation style, view size and vault state on the right, as text. & \eref{UIX-06}\\
\bottomrule
\end{tabularx}
\end{center}

\subsubsection*{What ``clean but feature dense'' means in numbers}

Dense means more on screen without a smaller target. Clean means one of everything: one icon size per place,
one corner radius per kind of thing, one accent. Every size is a multiple of 2\,px on a 4\,px grid.

\begin{center}
\small
\begin{tabularx}{\linewidth}{@{}lXl@{}}
\toprule
Thing & Rule & Today\\
\midrule
Text & 12\,px system font (SF Pro); captions 10\,px semibold capitals, letter-spaced & 13\,px (\file{BruceWorx.qss})\\
Tree rows & 22\,px; property rows 21\,px; dialog fields 28\,px & tree 28\,px, task list 42\,px\\
Icons & 24\,px toolbar, 20\,px timeline and navigation, 16\,px tree and tabs & toolbar about 22 and 40\,px mixed\\
Corners & 6\,px controls, 10\,px cards & mixed\\
Panels & Floating cards over the 3D view with a translucent, blurred background & docked columns\\
Colour & Neutral greys plus \emph{one} accent; state colours only for state (WIP, released, errors) & several blues\\
\bottomrule
\end{tabularx}
\end{center}

With 22\,px rows the Browser shows $\lfloor 560/22 \rfloor = 25$ rows in the height that holds 20 today. The ``today'' column is measured from figure~\ref{fig:ui-before} and the stylesheet. The
command dialog fits the whole Revolute joint, limits included, without scrolling.

\subsubsection*{Light and dark, following macOS}

\begin{figure}[!htbp]
\centering
\includegraphics[width=0.497\linewidth]{fig/ui_mock_part-light.png}\hfill
\includegraphics[width=0.497\linewidth]{fig/ui_mock_part-dark.png}\\[3pt]
\includegraphics[width=0.497\linewidth]{fig/ui_mock_asm-light.png}\hfill
\includegraphics[width=0.497\linewidth]{fig/ui_mock_asm-dark.png}
\caption{Part modelling and assembly, light and dark. The assembly shows the Revolute joint from section
\ref{sec:joints}. The limits get a range bar with the current reading, ``Measured from'' shows both datums, and
the angle is labelled in the 3D view. Assemblies have no feature history, so the timeline row says so instead of
showing a PartDesign message.}
\label{fig:ui-part-asm}
\end{figure}

\begin{figure}[!htbp]
\centering
\includegraphics[width=0.497\linewidth]{fig/ui_mock_sketch-light.png}\hfill
\includegraphics[width=0.497\linewidth]{fig/ui_mock_sketch-dark.png}\\[3pt]
\includegraphics[width=0.497\linewidth]{fig/ui_mock_start-light.png}\hfill
\includegraphics[width=0.497\linewidth]{fig/ui_mock_start-dark.png}
\caption{Sketch and Start page, light and dark.
\emph{Sketch:} the contextual \textsc{Sketch} tab replaces the solid tools, with \textsc{Finish Sketch} at the end. The sketch palette shows the solver (``2 degrees of freedom'', with the free geometry in blue) and a dense constraint list. The timeline is rolled back to the sketch.
\emph{Start:} new and open on the left, with the vault summary; recent files as cards, then as a table with P/N, revision and state.}
\label{fig:ui-sketch-start}
\end{figure}

\paragraph{How following the system works.} FreeCAD 1.1 has no setting for this. It applies one theme
(\file{BruceWorx.yaml} or \file{BruceWorx Dark.yaml}) and one stylesheet at start-up. This build's Qt is 6.8.3,
and Qt reports the system appearance as \code{QStyleHints::colorScheme()}, with a \code{colorSchemeChanged} signal
when you switch. BruceWorx listens for that signal. It loads the matching yaml colour tokens, re-applies
\file{BruceWorx.qss} (already written once for both themes), and sets the two viewport background colours and the
NaviCube colours. No restart is needed. The accent comes from \code{QPalette::Accent} (Qt~6.6 and later), which
on macOS is the System Settings accent. It is written into the \code{@AccentColor} token, and the derived tokens
are computed from it:
\[
  \text{accent-soft} = \alpha\,A + (1-\alpha)\,B, \qquad \alpha = 0.14 \text{ (light)},\ 0.22 \text{ (dark)},
\]
where $A$ is the accent and $B$ the card colour. Selected rows, pressed tools and chips use accent-soft. Only the
default button, checkboxes and the focus ring use the full accent.

\paragraph{Contrast, checked rather than eyeballed.} With a user-chosen accent, text on it can become
unreadable (white on yellow). The rule uses the WCAG relative luminance of each colour,
\[
  L = 0.2126\,R + 0.7152\,G + 0.0722\,B, \qquad c = \begin{cases} c_s/12.92 & c_s \le 0.04045\\ \left(\tfrac{c_s+0.055}{1.055}\right)^{2.4} & \text{otherwise}\end{cases}
\]
for each sRGB channel $c_s \in [0,1]$, and the contrast ratio $C = (L_1+0.05)/(L_2+0.05)$ with $L_1 \ge L_2$.
Text on the accent is white when $C_\text{white} \ge 3$ (macOS's own choice for buttons), otherwise near-black.

\begin{center}
\small
\begin{tabular}{@{}lrr@{\qquad}lrrl@{}}
\toprule
Text on card & Light & Dark & macOS accent & White & Black & Text\\
\midrule
Primary & 16.8 & 11.8 & Blue & 4.02 & 5.23 & white\\
Secondary & 5.1 & 5.4 & Purple & 4.88 & 4.31 & white\\
Tertiary (hints) & 3.6 & 3.5 & Pink / Red & 3.20 / 3.18 & 6.6 & white\\
\textit{Faded tree today} & \textit{1.1} & --- & Orange & 2.57 & 8.16 & black\\
 & & & Yellow / Green & 1.58 / 2.43 & 13.3 / 8.6 & black\\
\bottomrule
\end{tabular}
\end{center}
The mock's tertiary grey was darkened from 2.6:1 to 3.6:1 by this check. The accent values are Apple's published
light-mode colours. \verify

\begin{center}
\small
\begin{tabular}{@{}lll@{\qquad}lll@{}}
\toprule
Token & Light & Dark & Token & Light & Dark\\
\midrule
Chrome & \code{\#E9E9EC} & \code{\#232326} & Text & \code{\#1D1D1F} & \code{\#ECECEF}\\
Toolbar & \code{\#F3F3F5} & \code{\#2A2A2D} & Text 2 & \code{\#6E6E73} & \code{\#A1A1A6}\\
Card & white 88\,\% & \code{\#262629} 90\,\% & Text 3 & \code{\#86868B} & \code{\#808085}\\
Field & \code{\#FFFFFF} & \code{\#1C1C1E} & Accent & system & system\\
Viewport top & \code{\#FBFBFC} & \code{\#4A4F57} & OK & \code{\#1F9D55} & \code{\#32D74B}\\
Viewport bottom & \code{\#DFE3E9} & \code{\#26292E} & WIP & \code{\#C27C00} & \code{\#FFB340}\\
\bottomrule
\end{tabular}
\end{center}

\subsubsection*{Before and after}

\begin{figure}[!htbp]
\centering
\includegraphics[width=0.497\linewidth]{fig/ui_before_joint.png}\hfill
\includegraphics[width=0.497\linewidth]{fig/ui_mock_asm-light.png}
\caption{The same joint, today (left) and in the mock (right).}
\label{fig:ui-before-after}
\end{figure}

\subsubsection*{Plan: three steps, each shown before and after}

Each step ends with the four screens captured from the real build in light and dark, next to these mocks, for you
to approve.

\begin{center}
\small
\begin{tabularx}{\linewidth}{@{}lXl@{}}
\toprule
Step & Work & Size\\
\midrule
\textbf{A. Stylesheet and settings} & Colour tokens and sizes above in \file{BruceWorx.qss} and the two yaml files.
  Cards for Browser and dialog. NaviCube settings. Fix \marker{2}\marker{4}\marker{5}\marker{7}\marker{8}\marker{9}\marker{10}
  of figure~\ref{fig:ui-before}. Ribbon off in the BruceWorx profile. No C++.
  & small\\
\textbf{B. Small code} & Follow appearance and accent (signal + token reload). Document tabs on top. Dialog
  buttons at the bottom. Navigation pill. One-row timeline, hidden without history. Status-bar readouts.
  & medium\\
\textbf{C. Structure} & Workspace menu and tabs (\eref{UIX-05}), captioned groups (\eref{UIX-06}), tabs in the
  title bar, search, Start page. Planned in section~\ref{sec:ui-stepC} (approved 22 Oct), built in
  section~\ref{sec:ui-stepC-built} (Browser badges moved to step B). & large\\
\bottomrule
\end{tabularx}
\end{center}

\paragraph{Your answers (19 Oct, second round).}
\begin{enumerate}[nosep]
  \item \textbf{Dark-mode icons:} a faint light plate behind each icon in dark mode only, shown first
        (figure~\ref{fig:ui-plates}). The dark figures above already use it at 10\,\%.
  \item \textbf{Search:} ``yes''. Read as the built-in \texttt{Cmd-K} field; to be confirmed.
  \item \textbf{Ribbon add-on:} my call. BruceWorx turns FreeCAD-Ribbon off in its own profile while the
        BruceWorx bar is on and leaves it installed, so turning the bar off brings Ribbon back.
  \item \textbf{Properties:} stay in the Browser card, which grows down to fit them.
  \item \textbf{Assembly timeline:} components and joints in the order they were added. FreeCAD numbers every
        object as it is created (\code{obj.ID}), so the order needs no new data. No playback buttons:
        an assembly has nothing to roll back.
  \item \textbf{Document tabs:} a thin row under the title bar is acceptable if the title bar is impractical.
\end{enumerate}

\begin{figure}[!htbp]
\centering
\includegraphics[width=\linewidth]{fig/ui_mock_plates.png}
\caption{Dark mode, the same FreeCAD icons with no plate and with 6, 10 and 14\,\% white plates. The plate lifts
the black-outlined icons (construct, sketch create) clearly. It does much less for the red constraint icons,
whose problem is red on dark rather than outline on dark. In Qt the toolbar plate is one stylesheet rule on the
tool buttons. The plate behind tree and list icons needs a small item delegate (step~B).}
\label{fig:ui-plates}
\end{figure}

\subsubsection*{Step A: done (19 Oct 2026)}

\begin{figure}[!htbp]
\centering
\includegraphics[width=\linewidth]{fig/ui_stepA_sheet.jpg}
\caption{The real build, captured on a copy of your settings: before (Ribbon on), and after step A in light
and dark. Rows: part with a fillet open, the hinge joint, a sketch, the Start page. The 3D view is rendered
separately and pasted in, because a window grab leaves it blank; the view cube is not in these pictures.}
\label{fig:ui-stepA}
\end{figure}

\begin{center}
\small
\begin{tabularx}{\linewidth}{@{}X>{\raggedright\arraybackslash}p{5.1cm}@{}}
\toprule
Change & Where\\
\midrule
Colour tokens from the mock, light and dark; \code{blend()} gives the soft accent & \file{parameters/BruceWorx*.yaml}\\
12\,px type, 22\,px rows, 6/10\,px corners, flat task boxes as cards, thin tab $\times$ & \file{BruceWorx.qss}, \file{close-thin-*.svg}\\
Tree names no longer cut short (horizontal item padding removed) & \file{BruceWorx.qss}\\
Checkboxes show a tick; drop-down tools no longer paint a black bar & \file{BruceWorx.qss}\\
Dark-mode icon plate, 6\,\% white mixed into the toolbar colour (tokens carry no alpha) & \file{BruceWorx Dark.yaml}\\
Browser slightly transparent (your choice): a tinted glass card, 72\,\% opaque, text with a halo & \file{BruceWorx.py}, \file{BruceWorxModelTree.qss}\\
Task panel floats right as one card, shown only while a command is open; its boxes are flat sections & \file{BruceWorx.py}, \file{BruceWorxModelTree.qss}\\
OK/Cancel at the foot of every task panel (dialogs that choose a side keep it) & \file{TaskDialog.cpp}\\
Picks shown as lists: limit datums (Delete removes one), pattern objects & \file{JointObject.py}, \file{bruceworx\_pattern\_gui.py}\\
FreeCAD-Ribbon not loaded while \emph{BruceWorx top bar} is on (stays installed; stock FreeCAD still loads it) & \file{FreeCADInit.py}, \file{FreeCADGuiInit.py}\\
Two toolbar rows after every workbench switch (stale line breaks gave three or four) & \file{bruceworx\_toolbars.py}\\
Timeline: one row of icon chips in theme colours, no dock title, quiet empty message & HistoryTimeline mod\\
Start page readable in dark & \file{BruceWorx.qss}\\
Viewport gradient, accent, selection and view-cube colours from the mock & \file{bruceworx\_pack.py} (cfg regenerated)\\
Sketch grid on, as in the mock (your pick): faint translucent lines (white on dark, black on light), a major line every 5 squares, follows the zoom. The grid now honours a colour's alpha & \file{bruceworx\_pack.py}, \file{ViewProvider\-Grid\-Extension.cpp}\\
\emph{Sketch grid strength} preference, 2--40\,\%, default 10; open sketches update at once. Python: \code{bruceworx\_grid.set\_strength(14)}; a \code{GRID 14} alias comes with the command line (\eref{AUT-05}) & \file{bruceworx\_grid.py}\\
Joint panel \emph{Clear} is a normal button & \file{TaskAssemblyCreateJoint.ui}\\
\bottomrule
\end{tabularx}
\end{center}

New switches on Preferences $\rightarrow$ BruceWorx, all on by default: \emph{BruceWorx top bar}, \emph{Two
toolbar rows}, \emph{Float the task panel}; \emph{Transparent background} is on again by default. OK/Cancel at the
bottom reads \code{BruceWorx/TaskButtonsAtBottom} (no checkbox yet). Tests: 10 BruceWorx, 19 pattern, 16 Assembly, 26 pack, all passing.
To see the new colours on your own profile the pack must be applied again (Preferences $\rightarrow$ General
$\rightarrow$ Preference packs $\rightarrow$ BruceWorx), because a pack writes its values once.

\paragraph{Found on the way.} Editing a \file{.ui} panel does not rebuild the module's \file{\_rc.py}: the build
only watches the \file{.qrc}, and Python loads panels from that file. That was the cause of the stale joint
panel earlier today too. Until the CMake rule is fixed, delete \file{build/release/\{src/Mod/X/Gui,Mod/X\}/X\_rc.py}
after changing a panel.

\paragraph{Left for step B.} Task card header with the command's icon and title (a box without a title still
shows an empty strip); an inset from the right window edge (the overlay pins panels to the edge);
\textsc{Browser} header and badges; navigation pill; timeline in creation order for assemblies and hidden when
there is no history; follow the macOS appearance and accent live; sketch colours (free geometry blue,
constrained in text colour, dimensions quiet); document tabs on top; the property group captions.

\subsubsection*{Step B: built (21 Oct 2026), for your check}

\paragraph{1. Follow macOS appearance and accent (done).} \file{bruceworx\_appearance.py} listens for Qt's
\code{colorSchemeChanged} and for palette or activation events (macOS has no accent signal). On a light/dark
change it writes only the pack values that differ between the two modes; FreeCAD reloads the stylesheet when the
Theme parameter changes. The accent is read from \code{AppleAccentColor} and written to
\code{ThemeAccentColor1}, which the yaml files now use as \code{@AccentColor}, and to the 3D selection and
view-cube highlight. Text and ticks on the accent follow the contrast rule above: verified live, a yellow
accent gets dark text and dark ticks, blue gets white. Two switches on Preferences $\rightarrow$ BruceWorx, both on.
Theme tokens from \emph{Themes/UserParameters} take precedence over the yaml, so the computed
\code{AccentTextColor} wins over the yaml default.

\paragraph{2. Sketch colours (done).} In the pack, light and dark: free (under-constrained) geometry is blue
(\code{EditedEdgeColor}), a constrained element and a fully constrained sketch turn the text colour instead of
green, dimensions are a quiet text colour and constraint symbols the secondary text colour instead of red.
Construction lines are a lighter blue, grey once constrained; reference dimensions soft blue, dimensions driven by
an expression amber, external geometry purple. Sketches outside edit mode are drawn in a dark or light grey.
All in \file{bruceworx\_pack.py}.

\paragraph{3. Task card (done).} \file{bruceworx\_taskcard.py} gives the floating task panel a dialog header:
while a command is open the panel takes the title and icon of the dialog's first task box (``Fillet
Parameters'', ``Revolute joint''), read from the dialog itself, not from the workbench's command lists. That
box's own header is hidden so the title is not shown twice, and boxes with an empty header lose the empty strip;
later boxes keep their headers as section captions. When the command closes the panel is ``Tasks'' again with
every header back. The overlay title bar now draws the icon, then the title, from the left (\code{[BW]} in
\file{OverlayWidgets.cpp}). The panel has an 8\,px inset from the right window edge (\code{OverlayRight/Offset3}).
Only while \emph{Float the task panel} is on.

\paragraph{4. Timeline for assemblies, hidden without history (done).} For an assembly the timeline lists its
components and joints in the order they were created, sorted by \code{obj.ID}, as the same icon chips, with the
assembly's name on the left and ``5 items in the order they were added'' on the right. There is no marker and
nothing rolls back: a click selects the object in its assembly, a double-click on a joint opens it, the object
being edited is lit and a suppressed joint is struck through. The active body wins, then the active assembly,
then the document's assembly, then its first body with features. When there is nothing to show (no document, an
empty body or assembly) the bar hides and comes back as soon as there is. New switch on Preferences
$\rightarrow$ BruceWorx, on: \emph{Hide the history timeline when there is no history}
(\code{Mod/HistoryTimeline/HideWhenEmpty}). The PartDesign timeline is unchanged. Tests: 36 new headless checks
(\file{test\_assembly\_timeline.py}), the 71 existing ones pass.

\paragraph{5. Navigation pill (done).} \file{bruceworx\_navpill.py} puts a small card at the bottom centre of
every 3D view, as in the mock: orbit, pan, zoom and fit, then display style, camera, section and measure. The right
group runs FreeCAD's own commands (\code{Std\_DrawStyle} as a menu, orthographic/perspective in one button,
\code{Std\_ToggleClipPlane}, \code{Std\_Measure}) with their icons. Orbit, pan and zoom are modes as in Fusion:
click, drag in the view, Esc to stop; orbit turns about the vertical axis through the focal point. The pill moves
into the space between the Browser and the task card when the view is narrow, and closes with its view. Colours
come from a \code{\#BruceWorxNavPill} block in \file{BruceWorx.qss}; dark mode has the same faint icon plate as
the toolbar. Switch on Preferences $\rightarrow$ BruceWorx, on. Tests: 20 pure-logic tests; every button clicked
in a GUI run, light and dark captured.

\paragraph{5a. Navigation pill orbit centre (done).} The pill's orbit turns either about the point the camera looks
at (the focal point) or about the last clicked point: the point on the model under the pointer where the drag
starts, or the focal plane there when nothing is under it. The pick is the one a mouse orbit makes
(\code{NavigationStyle::saveCursorPosition}). It becomes FreeCAD's rotation centre, with ``Window center'' treated
as ``Drag at cursor''. A small [BW] binding exposes \code{getRotationCenter}, \code{setRotationCenter} and
\code{findRotationCenter} on \code{NavigationStyle} to Python. Choice on Preferences $\rightarrow$ BruceWorx
(``Navigation pill orbits about'', \code{NavPillOrbitCenter}), default the last clicked point. Tests: 5 pure-logic
tests for the pivoted orbit; a GUI run checked that the clicked model point stays fixed on screen in one mode and
the focal point in the other.

\paragraph{6. Document tabs on top and status-bar readouts (done).} The tab position is a preference,
\code{MainWindow/DocumentTabPos} (stock \code{South}; the pack sets \code{North}). The tabs sit in a thin row under
the toolbars, and the Browser and Tasks panels now start below them. A modified document shows a dot in place of
the close button, with the $\times$ on hover; unmodified tabs show it on the current tab only
(\code{BruceWorx/DocumentTabDot}). The status bar reads as text (\code{BruceWorx/StatusReadouts}): the prompt on the
left; on the right \emph{Report} (warnings in orange, errors in red; a click opens the Report view), \emph{Units} (the
unit menu), \emph{Navigation}, \emph{View} (3D views only) and \emph{Vault}, read from the folder vault's ref file named
by \code{BruceWorx/VaultPath} and hidden until one is set. The check-out state waits for a cached vault status.
Status messages no longer turn black on the dark bar. Three switches on Preferences $\rightarrow$ BruceWorx, on
through the pack; stock FreeCAD is unchanged. \file{MainWindow.cpp}, \file{OverlayManager.cpp},
\file{bruceworx\_statusbar.py}, \file{BruceWorx.qss}.

\paragraph{7. Browser header, badges and captions (done).} The left card is titled \textsc{Browser}. A part or assembly
with the part properties (\S\ref{sec:partprops}) shows its part number and state as badges after its name
(``LE-01-01001'', ``WIP 02'' in the warning colour); a component in an assembly shows them as quiet text. Features
get a quiet summary of their main dimension (``Pad 8 mm'', ``R5'', ``\O 9 $\times$ 4''); the type is named only when
the label does not say it. Badges and summaries follow every change. In dark mode a 10\,\% white plate sits behind
each tree icon. Property groups are quiet small-caps captions instead of grey bars. The state comes from a
\code{State} property when there is one, otherwise from the revision code (numeric = WIP, letter = released). The
tree draws the extras itself and asks \code{FreeCADGui.BrowserRowInfo} once per change (\file{Tree.cpp},
\file{bruceworx\_browser.py}, \file{PropertyItemDelegate.cpp}). Five switches on Preferences $\rightarrow$
BruceWorx, all on. Tests: 18 new.

\paragraph{Found while testing (20--21 Oct).}
\begin{itemize}[nosep]
  \item \textbf{Quit crash on macOS.} Quitting (Dock $\rightarrow$ Quit, logout, \code{QApplication::quit()})
        crashed and the settings were not saved: \code{MainWindow::closeEvent} called \code{quit()} again, the
        nested \code{[NSApp terminate:]} made AppKit call \code{exit()} at once, and the static destructors
        crashed. The quit is now queued until the close returns (macOS only, \file{MainWindow.cpp}).
  \item \textbf{No top toolbar.} FreeCAD-Ribbon hides FreeCAD's own toolbars and FreeCAD stores that per toolbar,
        so with the Ribbon replaced the window could start with no toolbar at all.
        \code{bruceworx\_toolbars.show\_if\_all\_hidden} runs after every workbench switch: when all app-wide
        toolbars, or all of the workbench's, are hidden it shows them again. A toolbar hidden on purpose stays
        hidden while another of its group shows.
  \item \textbf{Task watchers off.} The workbench command lists FreeCAD shows in the task panel (Insert,
        Constraints, \ldots) repeat the toolbar; the pack turns them off (\code{General/ShowTaskWatcher}), with a
        checkbox on Preferences $\rightarrow$ BruceWorx.
  \item \textbf{Trackpad gestures.} A plain two-finger swipe zooms as before; Shift + swipe pans in any direction,
        Cmd + swipe orbits about the point under the pointer. The view follows the fingers (natural scrolling
        honoured); a mouse wheel is not affected. \code{NavigationStyle::processTrackpadScroll}, called from the
        viewer's wheel filter; switch \code{BruceWorx/TrackpadGestures}, on.
  \item \textbf{2\,px lines.} \code{DefaultShapeLineWidth} is 2 again: Coin line widths are device pixels, so
        1\,px is a half-point hairline on Retina. Body edges and sketches outside edit mode.
  \item \textbf{Grid when a sketch opens.} FreeCAD stores ``show grid'' in each sketch when it is created, so
        older sketches opened without one. \file{bruceworx\_grid.py} switches it on whenever a sketch opens for
        editing; the sketch panel can still hide it. Switch \code{BruceWorx/SketchGridOnEdit}, on.
  \item \textbf{Pattern while another command is open.} Rectangular or Circular Pattern did nothing when another
        task dialog was open; it now says ``Another command is open in the task panel. Close it (OK or Cancel)
        first.''
  \item \textbf{Checkable sections.} A group box that can be switched on (``Direction~2'') showed no checkbox.
        \file{BruceWorx.qss} now styles \code{QGroupBox::indicator} like a checkbox, caption in the text colour.
  \item \textbf{Checklist files.} \file{scripts/checklist\_files.py} builds, in the GUI and with a saved view,
        the starting file for each group of steps of the test session (section~\ref{sec:tonight}). The files go
        to \file{\textasciitilde/Developer/BruceWorx-checklist/}.
  \item \textbf{Start-up traceback.} With a file on the command line the ground grid asked for the scene graph
        before \code{pivy.coin} was imported (``No SWIG wrapped library loaded''). \file{ground\_grid.py} imports
        it first; the grid now appears on that file too.
  \item \textbf{Stale \file{\_rc.py} fixed.} The rule from step A (``Found on the way'') is fixed: \code{PYSIDE\_WRAP\_RC}
        reads the files listed in the \file{.qrc} at configure time and rebuilds \file{X\_rc.py} when any of
        them changes; editing the \file{.qrc} re-runs CMake. No more deleting \file{\_rc.py} by hand.
  \item \textbf{Review fixes (21 Oct).} A code review of the items above found four things, now fixed:
        \code{Structure} counts as an app-wide toolbar, so the toolbar restore fires again; the trackpad mode is
        chosen on the first swipe event that moves (macOS starts a scroll as soon as the fingers rest, before
        Shift or Cmd is pressed), and off macOS a wheel event without a scroll phase is never a swipe
        (high-resolution mice); the task card follows boxes a dialog shows later (Sketcher's tool box) and keeps
        the translated ``Tasks'' title.
  \item \textbf{Empty task card stayed on screen.} With task watchers off and the task view alone in its panel, the
        card never hid: \emph{Auto task} only folds a dock that shares its panel, the task view counted hidden
        watchers as content, and closing a dialog sent no update. All three fixed (\code{[BW]} in
        \file{OverlayWidgets.cpp}, \file{OverlayManager.cpp}, \file{TaskView.cpp}); the card now shows only while
        a command is open.
  \item \textbf{Checklist steps 8--17 run by script first (21 Oct).} Steps 8, 9, 12, 13, 14 and 16 passed; the
        failures are fixed and pass on a re-run. Step 10: the BOM counts the visible elements of a link array
        (6 bolts, not 1; \code{[BW]} in \file{BomObject.cpp}). Step 11: independent copies no longer leave the
        old array's \code{Pattern\_i\ldots} objects behind. Step 12: the feature's body is activated before
        PartDesign's own dialog, so no ``Active Body Required'' prompt. Step 13: the preferences page has an
        icon, and the pattern buttons go away without a workbench switch. Step 14: ``Now at 0.0\,$^\circ$'', not
        $-0.0$. Step 15: a joint whose part or face no longer exists shows an error in the tree and takes the
        part back when it returns. Step 17: Material name and Mass follow a material change at once; FreeCAD's
        placeholder ``Default'' material counts as no material. BOM cells for a property a row does not have
        stay empty instead of ``N/A'' with a warning each.
\end{itemize}
